% !TEX root = ../SANER2027-main-labeling.tex
% Threats to Validity

\section{Threats to Validity}
\label{sec:threats}

\myparagraph{Internal Validity} Each method's best $k$ is the one with the highest macro $F_1$ on the test set, which favors the RAG variants. With a single $k{=}12$ for every model size, \frag\ trails fine-tuning by 1.2 points on average instead of 1.0 (\Cref{sec:rq3}). The filter's margin was calibrated on a validation split of the labeled issues, not on the test set, and is the same for all models and projects (\Cref{sec:approach-bragtag}). Prompts and hyperparameters affect fine-tuning results~\cite{heo2025study,aracena2025applying}. Ours follow prior work where it reports them (\Cref{sec:finetune}), and other choices could change the results of either approach. Invalid outputs count as errors, which penalizes the RAG variants most; the fallback analysis shows their effect (\Cref{sec:rq3}). The issues are public and may appear in the LLMs' pretraining data, which would affect zero-shot prompting and both approaches, since all use the same models.

\myparagraph{Construct Validity} Macro $F_1$ on a balanced three-label test set weighs the labels equally, and we report per-class precision and recall so that trade-offs between labels stay visible (\Cref{tab:results}). Some ground-truth labels are wrong, accounting for 9\% of the sampled errors (\Cref{sec:disc-failures}). Since all methods are scored against the same labels, this noise lowers every method's scores. LLM outputs can vary between runs, so we ran the experiments three times and report the average. The failure analysis covers 150 sampled errors of Qwen-32B, so its shares are estimates and may differ at other sizes. Times depend on the GPU, and at Qwen-14B, fine-tuning and the RAG variants ran on different GPUs (\Cref{sec:setup-hardware}).

\myparagraph{Conclusion Validity} We test differences between methods with paired bootstrap CIs that resample the test issues (\Cref{sec:setup-metrics}). The per-size comparisons are not corrected for multiple comparisons, so a single per-size result should be read with care; the headline rests on the average over the four sizes. The CIs resample issues, not projects, so they do not capture the variation across projects, which reaches about 10 points on single projects (\Cref{fig:per-project-diff}).

\myparagraph{External Validity} We study one model family (Qwen2.5-Instruct, four sizes), one embedding model, and one benchmark of eleven GitHub projects with balanced labels. Results may differ with imbalanced labels, other issue trackers, or other model families. The $11\times$ ratio follows from pooling eleven projects with 300 labeled issues each, and other numbers of projects or of labeled issues per project may change the gap between the approaches. Issue text also varies in writing style, vocabulary, and structure across projects~\cite{akhavan2026linkanchorautonomousllmbasedagent}, and our results vary by project as well (\Cref{fig:per-project-diff}).
